% SEGMENT OLP-0023-S01
% Part: sets-functions-relations
% Chapter: functions
% Section: functions-relations

\documentclass[../../../include/open-logic-section]{subfiles}

\begin{document}

\olfileid{sfr}{fun}{rel}

\olsection{தொடர்புகளாகச் சார்புகள்}

\begin{explain}
$A$ இன் !!{element}s[acc] $B$ இன் !!{element}s உடன் இணைக்கும் ஒரு சார்பு,
$A$, $B$ ஆகியவற்றுக்கு இடையில் ஒரு தொடர்பை வரையறுக்கிறது என்பது தெளிவு.
$f(x) = y$ ஆக இருக்கும்போது, அப்போதும் அப்போது மட்டுமே, $x$, $y$
ஆகியவற்றுக்கு இடையில் அமையும் தொடர்பே அது.
உண்மையில், எடுத்துக்காட்டாகக் கணிதத்தின் அடிப்படைக் கூறுகளின் வகைகளைக்
குறைக்க விரும்பினால், சார்பு $f$ ஐ இந்தத் தொடர்புடன், அதாவது ஜோடிகளின்
ஒரு கணத்துடன், \emph{ஒன்றாகக் கருதவும்} முடியும்.
இதனால் ஒரு கேள்வி எழுகிறது: எந்தத் தொடர்புகள் இவ்விதமாகச் சார்புகளை
வரையறுக்கின்றன?
\end{explain}

% SEGMENT OLP-0023-S02
\begin{defn}[சார்பின் கோட்டுரு]
$f\colon A \to B$ ஒரு சார்பாக இருக்கட்டும்.
$f$ இன் \emph{கோட்டுரு} என்பது பின்வருமாறு வரையறுக்கப்படும்
$R_f \subseteq A \times B$ என்ற தொடர்பாகும்:
\[
R_f = \Setabs{\tuple{x,y}}{f(x) = y}.
\]
\end{defn}

% SEGMENT OLP-0023-S03
\begin{explain}
உறுப்புசார் சமத்துவத்தின்படி, ஒரு சார்பின் கோட்டுரு தனித்ததாக
நிர்ணயிக்கப்படுகிறது. மேலும், கணங்களுக்கான உறுப்புசார் சமத்துவம்,
சார்புகளுக்கான மறைமுக மதிப்புசார் சமத்துவக் கொள்கையை உடனடியாக
நியாயப்படுத்துகிறது: $f$, $g$ ஆகியவற்றுக்கு ஒரே சார்பகமும்
துணைச்சார்பகமும் இருந்தால், எல்லா மதிப்புகளிலும் அவை ஒத்திருக்கும்போது
அவை ஒரே சார்பாகும்.

அதேபோல, ஒரு தொடர்பு “சார்பாக அமையும்” தன்மையைக் கொண்டிருந்தால்,
அது ஒரு சார்பின் கோட்டுருவாகும்.
\end{explain}

% SEGMENT OLP-0023-S04
\begin{prop}\ollabel{prop:graph-function}
$R \subseteq A \times B$ பின்வரும் நிபந்தனைகளை நிறைவேற்றுவதாக இருக்கட்டும்:
\begin{enumerate}
\item $Rxy$, $Rxz$ எனில் $y = z$; மேலும்
\item ஒவ்வொரு $x \in A$ க்கும், $\tuple{x,
y} \in R$ என அமையும் ஏதாவது ஓர் $y \in B$ உள்ளது.
\end{enumerate}
அப்போது $R$ என்பது, $Rxy$ ஆக இருக்கும்போது அப்போதும் அப்போது மட்டுமே
$f(x) = y$ என வரையறுக்கப்படும் $f\colon A \to B$ என்ற சார்பின் கோட்டுருவாகும்.
\end{prop}

% SEGMENT OLP-0023-S05
\begin{proof}
$Rxy$ என அமையும் ஒரு $y$ இருப்பதாகக் கொள்வோம்.
$Rxz$ என அமையும் வேறொரு $z \neq
y$ இருந்தால், $R$ மீதான நிபந்தனை மீறப்படும்.
எனவே $Rxy$ என அமையும் ஒரு $y$ இருந்தால், இந்த $y$ தனித்ததாகும்;
ஆகவே $f$ தெளிவாக வரையறுக்கப்பட்டுள்ளது. $R_f = R$ என்பது தெளிவு.
\end{proof}

% SEGMENT OLP-0023-S06
\begin{explain}
ஒவ்வொரு சார்பு $f\colon A \to B$ க்கும் ஒரு கோட்டுரு உள்ளது.
அதாவது, $f(x) = y$ என்பதால் வரையறுக்கப்படும்,
$A \times B$ இல் அடங்கிய ஒரு தொடர்பு உள்ளது.
மறுபுறம், \olref{prop:graph-function} இல் கொடுக்கப்பட்ட பண்புகளைக்
கொண்ட ஒவ்வொரு தொடர்பு $R
\subseteq A \times B$ உம் ஒரு சார்பு $f \colon A \to
B$ இன் கோட்டுருவாகும்.
சார்புகளுக்கும் அவற்றின் கோட்டுருக்களுக்கும் உள்ள இந்த நெருங்கிய
தொடர்பால், ஒரு சார்பை அதன் கோட்டுருவாகவே கருதலாம்.
வேறு சொற்களில், சார்புகளைச் சில குறிப்பிட்ட தொடர்புகளுடன், அதாவது
வரிசைத் தொகுதிகளின் சில குறிப்பிட்ட கணங்களுடன், ஒன்றாகக் கருதலாம்.
\oliflabeldef{sfr:rel:ref:sec}{ஆனால் இந்த “ஒன்றாகக் கருதுதல்” என்பதன்
நோக்கம் \olref[sfr][rel][ref]{sec} இல் உள்ளதைப் போன்றதே என்பதைக்
கவனிக்கவும். இது சார்புகளின் மெய்ப்பொருளியல் பற்றிய கூற்று அல்ல;
சார்புகளைச் சில கணங்களாகக் \emph{கருதுவது} வசதியாக உள்ளது என்ற
கருத்தே ஆகும். இதன் வசதிக்கான ஒரு காரணத்தை இப்போது காணலாம்.}{}
தொடர்புகள் மீது செய்த செயல்களைப் போன்ற செயல்களைச் சார்புகள் மீதும்
செய்வதைக் கருதலாம் (\olref[sfr][rel][ops]{sec} ஐக் காண்க).
குறிப்பாக:
\end{explain}

% SEGMENT OLP-0023-S07
\begin{defn}\ollabel{defn:funimage}
$f \colon A \to B$ ஒரு சார்பாகவும் $C\subseteq A$ ஆகவும் இருக்கட்டும்.

$f$ ஐ $C$ க்குக் \emph{கட்டுப்படுத்துவதால்} கிடைப்பது,
அனைத்து $x \in C$ க்கும் $(\funrestrictionto{f}{C})(x) = f(x)$
என வரையறுக்கப்படும் $\funrestrictionto{f}{C}\colon C \to B$
என்ற சார்பாகும். வேறு சொற்களில்,
$\funrestrictionto{f}{C} = \Setabs{\tuple{x, y} \in R_f}{x \in
C}$.

$f$ ஐ $C$ மீது \emph{செயல்படுத்துவதால்} கிடைப்பது
$\funimage{f}{C} =
\Setabs{f(x)}{x \in C}$ ஆகும். இதை $f$ இன் கீழ் $C$ இன்
\emph{பிம்பம்} என்றும் அழைக்கிறோம்.
\end{defn}

% SEGMENT OLP-0023-S08
\begin{explain}
இந்த வரையறைகளிலிருந்து, எந்தச் சார்பு $f$ க்கும்
$\ran{f} =
\funimage{f}{\dom{f}}$ என்பது பின்தொடர்கிறது.
\oliflabeldef{sfr:rel:ops:sec}{\olref[sfr][rel][ops]{sec} இல் உள்ள
வரையறைகளையும் சார்புகளைத் தொடர்புகளாகக் கருதுவதையும் எடுத்துக்கொண்டால்,
இந்தக் கருத்துகள் ஒப்பானவையாக அமைகின்றன.
ஆயினும் இங்குள்ள சார்புக் கட்டுப்படுத்தல் உள்ளீட்டை மட்டும்
கட்டுப்படுத்துகிறது; முந்தைய இருமத் தொடர்புக்கான கட்டுப்படுத்தல்
இரு ஆயங்களையும் கட்டுப்படுத்துகிறது.
ஆனால் வேறு இரண்டு செயல்களான நேர்மாறுகளுக்கும் சார்புப் பெருக்கங்களுக்கும்
சற்று கூடுதல் விளக்கம் தேவை. அதனை \olref[inv]{sec},
\olref[cmp]{sec} ஆகியவற்றில் தருவோம்.}{}
\end{explain}

\end{document}
